\begin{figure}[tb]
\centering
\begin{tikzpicture}[x=0.57cm,y=0.6cm,font=\footnotesize,>=Latex]
\node[anchor=west] at (0,3.1) {\textbf{Batch: one long processing call}};
\draw[->] (0,1.85) -- (13.8,1.85);
\draw[fill=accent!25,draw=accent] (0.2,1.85) rectangle (1.1,2.85);
\node[anchor=west] at (1.5,2.45) {window $+$ FFT $+$ output};
\fill[accent] (0.65,1.85) circle (2pt);
\node[anchor=west] at (1.5,2.02) {publish immediately};
\node[anchor=west] at (0,0.95) {\textbf{Distributed: resume on successive calls}};
\draw[->] (0,-0.2) -- (13.8,-0.2);
\foreach \x in {0.2,1.2,2.2} {
  \draw[fill=ink!20,draw=ink] (\x,-0.2) rectangle +(0.65,0.48);
}
\foreach \x in {3.2,4.2,5.2,6.2,7.2,8.2,9.2} {
  \draw[fill=ink!65,draw=ink] (\x,-0.2) rectangle +(0.65,0.48);
}
\foreach \x in {10.2,11.2,12.2} {
  \draw[fill=accent!45,draw=accent] (\x,-0.2) rectangle +(0.65,0.48);
}
\node[anchor=north,align=center] at (1.5,-0.3) {window /\\pack};
\node[anchor=north] at (6.5,-0.3) {FFT butterflies};
\node[anchor=north,align=center] at (11.5,-0.3) {magnitudes /\\output};
\draw[<->] (0.2,-1.5) -- (12.85,-1.5);
\node[fill=white,inner sep=2pt] at (6.5,-1.5) {one hop: $H$ sample calls};
\draw[->,accent] (12.85,0.85) -- (12.85,0.3);
\node[anchor=east] at (12.65,0.7) {publish};
\end{tikzpicture}
\caption{The same selected frame can be computed at its endpoint or spread
over the following hop. Input capture continues in both cases. Bars are
conceptual work slices, not timing measurements; stages and bar heights are
not to scale. The distributed result is published on the last call.}
\label{fig:schedule}
\end{figure}
